\section{Discussion}
\label{sec:discussion}

The discussion is most useful when three questions are separated: what explicit self-consistency adds beyond simpler charge-aware corrections in static fitting, how that increment behaves once trajectories leave the near-equilibrium manifold, and what the final CL-20/HMX target implies about the present scope of charge-aware stabilization.

We ablate the self-consistent charge-equilibration module by comparing full CSC-MACE against \texttt{static\_charge}, where the equivariant features predict a simpler non-self-consistent charge correction, and \texttt{short\_range}, which excludes explicit electrostatics entirely. On OpenREACT, the self-consistent branch achieves a mean force MAE of $3.25 \times 10^{-2} \pm 2.49 \times 10^{-3}$, improving on both alternatives, but this force gain is accompanied by a slight energy penalty relative to \texttt{static\_charge}.

That tradeoff is methodologically informative because it shows that explicit self-consistency is not redundant with a simpler charge-aware head. The persistent energy--force split indicates that self-consistency captures higher-order field response, while the present parameterization leaves room for tighter coupling between the electrostatic and short-range terms.

The Transition1x stress tests reveal the first clear limitation of that increment: static gains of the charge-aware models do not translate directly into enhanced trajectory stability under extreme shock protocols. All families synchronize at the onset of local bond-network rearrangement, but the \texttt{short\_range} baseline preserves a larger post-onset stability margin, whereas both charge-aware branches enter the degraded regime earlier. This suggests that the added electrostatic degrees of freedom, while physically useful near equilibrium, may introduce additional sensitivity in the highly distorted configurations characteristic of post-onset shock dynamics.

The real CL-20/HMX benchmark reveals a different and more decisive limitation. Static energy ranking improves monotonically from \texttt{short\_range} to \texttt{static\_charge} to \texttt{self\_consistent}, but that hierarchy does not produce a corresponding nonequilibrium separation under the current rollout protocols. Across stress, threshold-aware stress, reaction-progress trajectories, dense frame ensembles, and focused key-bond survival comparisons, the dominant variation is the starting frame rather than the model family.

In particular, molecular-integrity loss and C--N / N--O survival remain almost identical across the three families once the frame is fixed, while the same family's spread across frames is much larger. The focused dominance summary makes that point explicit: on the final key-bond analysis, the mean spread across families within a fixed frame is zero for the final molecular-integrity and retained-bond endpoints, whereas the frame-driven spread is about $8\times$ larger for the first integrity crossing and about $184\times$ larger for the first drift-threshold step.

A dedicated frame-$160$ nitro-focus multiseed analysis strengthens this limitation rather than weakening it: \texttt{short\_range} retains a mean nitro N--O fraction of $0.444$, whereas \texttt{static\_charge}, \texttt{self\_consistent}, and two damped self-consistent ablations all cluster at $0.463$; meanwhile, the nitro-motif integrity converges to the same mean value ($0.148$) across all five families, with first motif-integrity threshold crossings confined to the narrow band of steps $1647$--$1663$.

We then probed a stronger methodological intervention: a stability-gated self-consistent branch and an adaptive-gated branch that modulates the self-consistent mixing step size. On frame $160$, neither created a new family-level outcome beyond the simpler charge-aware plateau at step $1067$, and matched profiling shows a clear rollout-cost ordering on that same frame, from about $20$ s for \texttt{short\_range} to $42$ s for \texttt{static\_charge} to $112$ s for \texttt{self\_consistent}, before the adaptive-gated branch rises to about $190$ s while preserving the same key structure-survival outcome.

The practical boundary is therefore stronger than a simple reversal of static gain. On the final target, current CSC-MACE variants do not convert their static energy advantage into a reliably distinguishable dynamical event hierarchy, and more elaborate charge-aware updates can substantially increase computational burden without yielding a new structure-level advantage.

We emphasize that the effective charges predicted by CSC-MACE are parameters of a self-consistent interatomic potential and should not be over-interpreted as unique physical observables equivalent to any single DFT population analysis scheme such as Mulliken or Bader. Their primary role here is to improve force-field fidelity through the inclusion of long-range electrostatics. Recent charge-aware architectures have reached similar conclusions in adjacent settings: explicit charge degrees of freedom can be made active, interpretable, and useful, while the transfer of that information to non-equilibrium stability remains system-dependent \cite{fuchs2025celli,jung2026charge}. Whether closer alignment with explicit electronic-density observables would stabilize long-horizon dynamics remains an open question.

Taken together, these results argue for a layered validation strategy in which static error, early structural onset, fragment integrity, and post-onset stability are treated as distinct targets rather than interchangeable summaries. For energetic molecular crystals, the present CSC-MACE variants improve static energetics on the real target, but they do not induce a corresponding family-level separation in the tested nonequilibrium event metrics.

Even stronger adaptive interventions tested directly on the most fragile CL-20/HMX frame only recover the existing charge-aware drift-threshold plateau while increasing rollout cost, showing that more adaptive self-consistency alone does not create a new dynamical hierarchy. The paper therefore defines a method boundary with unusual specificity: explicit self-consistency is active, measurable, and beneficial for the real-target static hierarchy, but its extra physical structure does not dominate the present nonequilibrium outcome.

The next challenge is therefore sharply defined: stabilize the electrostatic gain so that it propagates into dynamically relevant observables without disproportionate computational overhead. The current data boundary is equally clear: the CL-20/HMX release already supports a real force-labeled energetic-material benchmark, and a held-out HMX-only split table can now be regenerated from the GPU multiseed static analysis, but the present contiguous split still does not provide simultaneous held-out CL-20 and HMX coverage or the reviewer-complete thermodynamic and labeling provenance needed for a final pressure- or trajectory-history-resolved mechanistic decomposition.
